% Research paper template built on ulpaper.cls.
%
% Layout: twocolumn (ICML-like, default) or onecolumn (NeurIPS-like).
% Add a4paper for A4, review for line numbers, numbers for [1]-style citations.
% There is no page limit; long papers put supporting material in appendices.
\documentclass[twocolumn]{ulpaper}
% \documentclass[onecolumn]{ulpaper}

% Optional packages for this paper
\usepackage{tikz}
\usetikzlibrary{arrows.meta,positioning,calc}
\usepackage{siunitx}

% Paper-specific macros
\newcommand{\method}{\textsc{Method}}
\DeclareMathOperator*{\argmax}{arg\,max}

\hypersetup{
  pdftitle={Paper Title: A Descriptive Subtitle},
  pdfauthor={First Author, Second Author},
  pdfkeywords={keyword one, keyword two},
}

\title{Paper Title:\\A Descriptive Subtitle}
\shorttitle{Paper Title}
\author{First Author\affil{1} \And Second Author\affil{1,2} \AND Third Author\affil{2}}
\affiliations{\affil{1}First Institution, City, Country.
  \affil{2}Second Institution, City, Country.}
\correspondence{First Author \textless\texttt{first@example.org}\textgreater.}
\notice{Preprint.}
\keywords{keyword one; keyword two; keyword three}

\begin{document}

\maketitle

\begin{abstract}
State the problem in one or two sentences. Say what is missing from current approaches. Describe the method or argument in two or three sentences. Report the main result with a number where there is one. Close with what the result means for the field. Aim for 150 to 250 words.
\end{abstract}

\section{Introduction}
\label{sec:intro}

Open with the problem and why it matters. Cite prior work in author-year form: \citet{vaswani2017attention} for a citation that is part of the sentence, and \citep{kingma2015adam} for a parenthetical one.

The second paragraph says what current methods do and where they fall short. The third paragraph introduces the approach of this paper.

We make three contributions.
\begin{enumerate}
  \item A formulation of the problem that makes the assumptions explicit (\cref{sec:method}).
  \item A method that addresses it, given in \cref{alg:method}.
  \item Experiments that test the method against baselines (\cref{sec:experiments}).
\end{enumerate}

\section{Related Work}
\label{sec:related}

\paragraph{First line of work.} Group related papers by idea and say how this paper differs from each group.

\paragraph{Second line of work.} Use run-in bold paragraph heads for each group instead of bullet lists.

\section{Method}
\label{sec:method}

\subsection{Problem formulation}

Define notation once, then use it consistently. Number only the equations that the text refers to.
\begin{equation}
  \label{eq:objective}
  \mathcal{L}(\theta) = \frac{1}{N}\sum_{i=1}^{N} \ell\bigl(f_\theta(x_i), y_i\bigr) + \lambda \lVert \theta \rVert_2^2 .
\end{equation}
\noindent Text that continues the sentence after a display starts with \verb|\noindent|.

\begin{assumption}[Bounded loss]
\label{asm:bounded}
The loss satisfies $0 \le \ell \le 1$.
\end{assumption}

\begin{definition}[Score]
\label{def:score}
For an input $x$, the score is $s(x) = \max_{y} p_\theta(y \mid x)$.
\end{definition}

\subsection{Algorithm}

\Cref{alg:method} gives the procedure and \cref{fig:overview} shows its components.

\begin{algorithm}[t]
\caption{Training with \method}
\label{alg:method}
\begin{algorithmic}[1]
\Require data $\{(x_i, y_i)\}_{i=1}^N$, step size $\eta$, steps $T$
\State initialize $\theta_0$
\For{$t = 1, \dots, T$}
  \State sample a minibatch $B$
  \State $\theta_t \gets \theta_{t-1} - \eta \nabla \mathcal{L}_B(\theta_{t-1})$
\EndFor
\State \Return $\theta_T$
\end{algorithmic}
\end{algorithm}

\begin{figure}[t]
\centering
\begin{tikzpicture}[font=\small,
  box/.style={draw, rounded corners=2pt, fill=#1, minimum height=6mm, inner xsep=4pt, align=center},
  box/.default=white, arr/.style={-{Stealth[length=2mm]}, thick}]
  \node[box=ulblue!12] (a) {Input};
  \node[box, right=6mm of a] (b) {Encoder};
  \node[box=ulorange!20, right=6mm of b] (c) {Output};
  \draw[arr] (a) -- (b);
  \draw[arr] (b) -- (c);
\end{tikzpicture}
\caption{A column-width figure. The caption says what the reader should notice, in full sentences.}
\label{fig:overview}
\end{figure}

\section{Experiments}
\label{sec:experiments}

\subsection{Setup}

\begin{table*}[t]
\caption{A full-width table for content too wide for one column.}
\label{tab:wide}
\centering
\small
\begin{tabularx}{\textwidth}{@{}l X c c c@{}}
\toprule
Setting & Description & Metric 1 & Metric 2 & Metric 3 \\
\midrule
A & A setting that needs a longer description than a narrow column allows. & 0.81 & 0.77 & 0.92 \\
B & A second setting described at similar length. & 0.84 & 0.79 & 0.90 \\
\bottomrule
\end{tabularx}
\end{table*}

Describe data, baselines, metrics, and compute. \Cref{tab:main} reports the main results. Tables use \texttt{booktabs} rules only, with captions above.

\begin{table}[t]
\caption{Main results. Mean and standard deviation over three seeds; best in bold.}
\label{tab:main}
\centering
\small
\begin{tabular}{@{}lcc@{}}
\toprule
Method & Accuracy (\%) & Time (s) \\
\midrule
Baseline A & $85.2 \pm 0.4$ & 12.3 \\
Baseline B & $87.5 \pm 0.6$ & 15.7 \\
\method{} (ours) & $\mathbf{91.3} \pm 0.3$ & 14.2 \\
\bottomrule
\end{tabular}
\end{table}

\subsection{Results}

Report what the numbers show, then what they mean. Wide content that spans both columns uses the starred environments, as in \cref{tab:wide}. Place a starred float in the source before the page where it should appear, since LaTeX sets it at the top of the next page; in the one-column layout the starred forms behave like the plain ones.


\section{Discussion and Limitations}
\label{sec:limitations}

State what the results do not show, which assumptions may fail, and what remains open.

\section{Conclusion}
\label{sec:conclusion}

Summarize the contribution and the main result in a short paragraph, then name the next steps.

\section*{Acknowledgments}
Thank people and funders here.

\bibliography{references}

\appendix

\section{Additional Details}
\label{app:details}

Appendices hold proofs, extra results, hyperparameters, and code. They start on a new page and are lettered.

\begin{lstlisting}[language=Python, caption={Code listings use the class style.}, label={lst:example}]
def score(probs):
    """Return the maximum class probability."""
    return max(probs)
\end{lstlisting}

\end{document}
